\section{Certified localization and a usable numerical budget}
\label{sec:cert}
\subsection{Exact finite Euler evaluation on the axis}
To turn the characterization of $\Cs$ into a short rational enclosure, we use only finite Euler sums and the just-proved remainder bound. Rearranging \eqref{eq:Euler} gives
\begin{equation}\label{eq:finiteweights}
 E_N=2^{-N}\sum_{k=0}^{N-1}(-1)^k A_k
                  \sum_{j=k+1}^{N}\binom Nj.
\end{equation}
The finite coefficient identity behind this formula is
\[
 \sum_{j=k}^{N-1}2^{N-j-1}\binom jk
 =\sum_{j=k+1}^{N}\binom Nj.
\]
It follows by induction from Pascal's identity. The supplied check compares the two implementations for all $1\le N\le32$, or $528$ individual coefficients.

On the axis $a=0$, $A_k=H_{2k}^{(b)}$ and $C(b)=-2g_{0,b}$. Suppose exact rational arithmetic encloses the finite sum by $L\le E_N\le U$. For $N\ge2$, the kernel theorem implies
\begin{equation}\label{eq:axisinterval}
 -2U<C(b)<-2L+\frac{9}{4}\,2^{-N}.
\end{equation}
This is an enclosure of an analytic value, even if the original alternating series diverges.

For a positive rational order $b=p/q$ in lowest terms, each $m^{-b}$ is enclosed on a grid of denominator $D$ by an integer $r$ satisfying
\begin{equation}\label{eq:rootcert}
 r^q m^p\le D^q<(r+1)^q m^p.
\end{equation}
These integer inequalities prove $r/D\le m^{-b}<(r+1)/D$. All powers, harmonic sums, signs and Euler weights are then combined with rational outward bounds. A decimal computation may propose $r$ when generating a certificate; the verifier accepts it only after checking \eqref{eq:rootcert} with integers. It never calls a transcendental evaluator.

We use $N=48$, $D=10^{30}$, and the three orders $1.3021$, $1.3022$, $1.3023$. There are $94$ power brackets per order, hence $282$ in total. The following displayed intervals are deliberately rounded outward more coarsely than the stored rational endpoints:
\begin{center}
\begin{tabular}{ccc}
\toprule
$b$ & Lower bound for $C(b)$ & Upper bound for $C(b)$\\
\midrule
$1.3021$ & $1.1365611027504$ & $1.1365611027505$\\
$1.3022$ & $1.1365611033275$ & $1.1365611033277$\\
$1.3023$ & $1.1365611030380$ & $1.1365611030381$\\
\bottomrule
\end{tabular}
\end{center}
The middle lower endpoint exceeds both outer upper endpoints. Strict unimodality therefore gives
\begin{equation}\label{eq:bstarcert}
 1.3021<b_*<1.3023.
\end{equation}
For example, $b_*\le1.3021$ would put all three points on the decreasing side and contradict the first comparison. The other exclusion is analogous.

\subsection{A curvature bound closes the global enclosure}
\begin{lemma}[A sufficient explicit curvature bound]\label{lem:curvature}
For $1\le b\le2$, $|C''(b)|<1$.
\end{lemma}
\begin{proof}
For a gamma variable $T$ of shape $b$, set $L=\log T-\psi(b)$. Differentiation of its density gives
\[
 \E L=0,\qquad \E L^2=\psi_1(b),\qquad
 C''(b)=\E\bigl[(\phi(T)-1)\{L^2-\psi_1(b)\}\bigr].
\]
The subtraction of one is legitimate because the density's second derivative integrates to zero. Since $0<\phi-1<(\sqrt2-1)/2<1/4$ and
\[
 \psi_1(b)=\sum_{n\ge0}(b+n)^{-2}\le\zeta(2)<2
 \qquad (b\ge1),
\]
we obtain $|C''(b)|<\tfrac14(\E L^2+\psi_1(b))<1$. The trigamma identity is classical \cite{DLMFpsi}. The elementary upper bound $\zeta(2)<2$ also follows by comparing its terms after the first with $\int_1^\infty x^{-2}\dd x$.
\end{proof}
On an interval of length $h$, a function whose second derivative has absolute value at most one lies at most $h^2/8$ above the linear interpolation of its endpoints. This follows from the interpolation remainder or by comparison with the parabola $(x-x_0)(x_1-x)/2$.
Apply this on the two intervals between the three certified orders, with $h=10^{-4}$. Together with \eqref{eq:bstarcert}, this bounds the maximum by the largest stored endpoint upper bound plus $1/(8\cdot10^8)$. Exact rational comparison yields
\begin{theorem}[Certified universal constant]\label{thm:certified}
The sharp universal Euler constant satisfies
\begin{equation}\label{eq:Cstarcert}
 \boxed{1.1365611033<\Cs<1.1365611046.}
\end{equation}
Consequently $1.136562$ is a valid rational universal constant for all $a,b>0$ and $N\ge1$. The simpler constant $57/50$ is also valid on the whole positive quadrant, not only on the critical triangle.
\end{theorem}
The interval in \eqref{eq:Cstarcert} is deliberately much shorter to state than the exact intermediate rational endpoints. Its width is $1.3\times10^{-9}$. It is a mathematical enclosure conditional only on the preceding ordinary proofs and finite integer checks, not a confidence interval from a numerical optimizer.

A separate $100$-digit evaluation gives the diagnostics
\begin{align*}
 b_*&\approx1.3022165871012412095923717015551655,\\
 \Cs&\approx1.1365611033395095609525863757799424,\\
 C''(b_*)&\approx-0.0866705616486695948236951284473855.
\end{align*}
Only \eqref{eq:bstarcert} and \eqref{eq:Cstarcert} are certified here. The longer decimals are retained as reproducible diagnostics, not silently promoted to rigorous digits.

\subsection{A practical stopping rule}
An exact or outward-rounded evaluation of $E_N$ immediately gives
\[
 g_{a,b}\in[L-1.136562\,2^{-N},\ U]\qquad(N\ge1)
\]
with the sharper lower endpoint $L-(9/8)2^{-N}$ whenever $N\ge2$. For a target remainder $\varepsilon>0$ and $N\ge2$, it suffices to choose
\[
 N\ge\left\lceil\log_2\frac{9}{8\varepsilon}\right\rceil.
\]
Any uncertainty in forming the finite sum must be included separately in $U-L$. In particular, this is not a license to ignore cancellation in floating-point binomial sums. The exact-weight formula \eqref{eq:finiteweights} and interval arithmetic provide a direct way to account for it.
